% Ilustrasi SubBytes: satu byte matriks state dipetakan melalui S-box menjadi
% byte baru. Digambar ulang sebagai diagram vektor; penggambaran aslinya dari
% slide Munir ditangkap pada bingkai animasi yang sorotannya meleset sehingga
% kotak sorot menutupi sel tetangga dan panahnya menunjuk baris yang salah.
% Dirujuk dari section-02.tex dengan % Ilustrasi SubBytes: satu byte matriks state dipetakan melalui S-box menjadi
% byte baru. Digambar ulang sebagai diagram vektor; penggambaran aslinya dari
% slide Munir ditangkap pada bingkai animasi yang sorotannya meleset sehingga
% kotak sorot menutupi sel tetangga dan panahnya menunjuk baris yang salah.
% Dirujuk dari section-02.tex dengan % Ilustrasi SubBytes: satu byte matriks state dipetakan melalui S-box menjadi
% byte baru. Digambar ulang sebagai diagram vektor; penggambaran aslinya dari
% slide Munir ditangkap pada bingkai animasi yang sorotannya meleset sehingga
% kotak sorot menutupi sel tetangga dan panahnya menunjuk baris yang salah.
% Dirujuk dari section-02.tex dengan % Ilustrasi SubBytes: satu byte matriks state dipetakan melalui S-box menjadi
% byte baru. Digambar ulang sebagai diagram vektor; penggambaran aslinya dari
% slide Munir ditangkap pada bingkai animasi yang sorotannya meleset sehingga
% kotak sorot menutupi sel tetangga dan panahnya menunjuk baris yang salah.
% Dirujuk dari section-02.tex dengan \input{figures/subbytes-aes}.
\begin{figure}[htbp]
    \centering
    \begin{tikzpicture}[
        sel/.style={draw, minimum size=1.15cm, inner sep=0pt, font=\small},
        sorot/.style={sel, fill=blue!55, text=white},
        panah/.style={-{Latex[length=2.4mm]}, thick},
    ]
        % Grid state (kiri) dan state' (kanan). Indeks mengikuti penamaan slide:
        % baris pertama a_00..a_03, dan seterusnya.
        \foreach \i in {0,1,2,3} {
            \foreach \j in {0,1,2,3} {
                \node[sel] (a\i\j) at ({\j*1.15}, {-\i*1.15}) {$a_{\i\j}$};
                \node[sel] (b\i\j) at ({\j*1.15+8.05}, {-\i*1.15}) {$b_{\i\j}$};
            }
        }
        % Byte yang sedang disubstitusi ditandai. Dipilih byte pada baris
        % terakhir agar panahnya dapat keluar tepat di tepi bawah grid tanpa
        % melintasi sel atau label lain.
        \node[sorot] at (a31) {$a_{31}$};
        \node[sorot] at (b31) {$b_{31}$};

        \node[font=\small] at (1.725, 1.15) {state};
        \node[font=\small] at (9.775, 1.15) {state$'$};

        \node[draw, thick, fill=white, minimum width=1.9cm, minimum height=1.1cm,
              font=\small] (sbox) at (5.75,-5.3) {S-box};

        % Byte tersorot masuk ke S-box, lalu keluar sebagai byte hasil substitusi.
        \draw[panah] (a31.south) to[out=-90, in=175] (sbox.west);
        \draw[panah] (sbox.east) to[out=5, in=-90] (b31.south);
    \end{tikzpicture}
    \caption{Transformasi SubBytes pada AES: setiap byte matriks \textit{state} disubstitusi melalui tabel S-box menjadi byte baru}
    \label{fig:subbytes-aes}
    \sumber{Munir, \textit{IF4020 Kriptografi}, ITB.}
\end{figure}
.
\begin{figure}[htbp]
    \centering
    \begin{tikzpicture}[
        sel/.style={draw, minimum size=1.15cm, inner sep=0pt, font=\small},
        sorot/.style={sel, fill=blue!55, text=white},
        panah/.style={-{Latex[length=2.4mm]}, thick},
    ]
        % Grid state (kiri) dan state' (kanan). Indeks mengikuti penamaan slide:
        % baris pertama a_00..a_03, dan seterusnya.
        \foreach \i in {0,1,2,3} {
            \foreach \j in {0,1,2,3} {
                \node[sel] (a\i\j) at ({\j*1.15}, {-\i*1.15}) {$a_{\i\j}$};
                \node[sel] (b\i\j) at ({\j*1.15+8.05}, {-\i*1.15}) {$b_{\i\j}$};
            }
        }
        % Byte yang sedang disubstitusi ditandai. Dipilih byte pada baris
        % terakhir agar panahnya dapat keluar tepat di tepi bawah grid tanpa
        % melintasi sel atau label lain.
        \node[sorot] at (a31) {$a_{31}$};
        \node[sorot] at (b31) {$b_{31}$};

        \node[font=\small] at (1.725, 1.15) {state};
        \node[font=\small] at (9.775, 1.15) {state$'$};

        \node[draw, thick, fill=white, minimum width=1.9cm, minimum height=1.1cm,
              font=\small] (sbox) at (5.75,-5.3) {S-box};

        % Byte tersorot masuk ke S-box, lalu keluar sebagai byte hasil substitusi.
        \draw[panah] (a31.south) to[out=-90, in=175] (sbox.west);
        \draw[panah] (sbox.east) to[out=5, in=-90] (b31.south);
    \end{tikzpicture}
    \caption{Transformasi SubBytes pada AES: setiap byte matriks \textit{state} disubstitusi melalui tabel S-box menjadi byte baru}
    \label{fig:subbytes-aes}
    \sumber{Munir, \textit{IF4020 Kriptografi}, ITB.}
\end{figure}
.
\begin{figure}[htbp]
    \centering
    \begin{tikzpicture}[
        sel/.style={draw, minimum size=1.15cm, inner sep=0pt, font=\small},
        sorot/.style={sel, fill=blue!55, text=white},
        panah/.style={-{Latex[length=2.4mm]}, thick},
    ]
        % Grid state (kiri) dan state' (kanan). Indeks mengikuti penamaan slide:
        % baris pertama a_00..a_03, dan seterusnya.
        \foreach \i in {0,1,2,3} {
            \foreach \j in {0,1,2,3} {
                \node[sel] (a\i\j) at ({\j*1.15}, {-\i*1.15}) {$a_{\i\j}$};
                \node[sel] (b\i\j) at ({\j*1.15+8.05}, {-\i*1.15}) {$b_{\i\j}$};
            }
        }
        % Byte yang sedang disubstitusi ditandai. Dipilih byte pada baris
        % terakhir agar panahnya dapat keluar tepat di tepi bawah grid tanpa
        % melintasi sel atau label lain.
        \node[sorot] at (a31) {$a_{31}$};
        \node[sorot] at (b31) {$b_{31}$};

        \node[font=\small] at (1.725, 1.15) {state};
        \node[font=\small] at (9.775, 1.15) {state$'$};

        \node[draw, thick, fill=white, minimum width=1.9cm, minimum height=1.1cm,
              font=\small] (sbox) at (5.75,-5.3) {S-box};

        % Byte tersorot masuk ke S-box, lalu keluar sebagai byte hasil substitusi.
        \draw[panah] (a31.south) to[out=-90, in=175] (sbox.west);
        \draw[panah] (sbox.east) to[out=5, in=-90] (b31.south);
    \end{tikzpicture}
    \caption{Transformasi SubBytes pada AES: setiap byte matriks \textit{state} disubstitusi melalui tabel S-box menjadi byte baru}
    \label{fig:subbytes-aes}
    \sumber{Munir, \textit{IF4020 Kriptografi}, ITB.}
\end{figure}
.
\begin{figure}[htbp]
    \centering
    \begin{tikzpicture}[
        sel/.style={draw, minimum size=1.15cm, inner sep=0pt, font=\small},
        sorot/.style={sel, fill=blue!55, text=white},
        panah/.style={-{Latex[length=2.4mm]}, thick},
    ]
        % Grid state (kiri) dan state' (kanan). Indeks mengikuti penamaan slide:
        % baris pertama a_00..a_03, dan seterusnya.
        \foreach \i in {0,1,2,3} {
            \foreach \j in {0,1,2,3} {
                \node[sel] (a\i\j) at ({\j*1.15}, {-\i*1.15}) {$a_{\i\j}$};
                \node[sel] (b\i\j) at ({\j*1.15+8.05}, {-\i*1.15}) {$b_{\i\j}$};
            }
        }
        % Byte yang sedang disubstitusi ditandai. Dipilih byte pada baris
        % terakhir agar panahnya dapat keluar tepat di tepi bawah grid tanpa
        % melintasi sel atau label lain.
        \node[sorot] at (a31) {$a_{31}$};
        \node[sorot] at (b31) {$b_{31}$};

        \node[font=\small] at (1.725, 1.15) {state};
        \node[font=\small] at (9.775, 1.15) {state$'$};

        \node[draw, thick, fill=white, minimum width=1.9cm, minimum height=1.1cm,
              font=\small] (sbox) at (5.75,-5.3) {S-box};

        % Byte tersorot masuk ke S-box, lalu keluar sebagai byte hasil substitusi.
        \draw[panah] (a31.south) to[out=-90, in=175] (sbox.west);
        \draw[panah] (sbox.east) to[out=5, in=-90] (b31.south);
    \end{tikzpicture}
    \caption{Transformasi SubBytes pada AES: setiap byte matriks \textit{state} disubstitusi melalui tabel S-box menjadi byte baru}
    \label{fig:subbytes-aes}
    \sumber{Munir, \textit{IF4020 Kriptografi}, ITB.}
\end{figure}
